\section{问题分析}
\subsection{总体分析}
五个问题不是孤立求值，而是沿着“底层资源--中层鱼群--珍稀物种--复杂食物链--综合情景”这条生态传导链展开。本文既要描述禁渔带来的恢复趋势，也要解释新失衡现象，并回答“为什么会恢复”“为什么会失衡”“何时会转向混沌”“污染和入侵会怎样放大脆弱性”。

\subsection{问题一分析}
问题一并不是简单比较禁渔前后的数量涨跌，而是要说明禁渔如何改变“底层资源--中层鱼类--上层捕食压力”之间的反馈方向。四大家鱼的食性不同，决定了它们对水草、浮游植物、浮游动物和底栖饵料的作用并不一致；因此，模型不能只给出末值比率，而应追踪恢复过程中的先后顺序、滞后关系和局地压力变化。这样才能解释同样处在禁渔背景下，为什么有些资源先回升，有些资源却会先被进一步挤压。

更具体地说，禁渔首先释放的是捕捞压力，随后才逐步改变种群之间的摄食关系。中层鱼类往往比底层资源更快恢复，这会带来更强的摄食和竞争效应，使局部水域出现“鱼类上升、底层下降”的阶段性现象。若只看末值，就会把“修复”和“失衡”混为一谈；只有把资源剩余、食压变化和恢复速度同时纳入，才能识别禁渔究竟是在推动系统回到更稳定的平衡，还是在形成新的局部承压。问题一的判据因此不应只是数量对比，还应包含能够反映资源压力和食压变化的综合指标。

从解释口径上看，这一问更适合按阶段来读：禁渔初期是捕捞压力迅速释放，鱼类数量先抬升；中期是中层鱼类对资源的摄食和空间竞争开始显现，底层资源的恢复被推迟；后期则要看系统能否回到新的平衡，而不是停留在局地高食压状态。也因此，问题一的结果应当用“资源压力是否回落、食压是否抬升、二者是否同步”来判断，而不是只写一个增长倍数。这样的分析既能解释底层资源为什么可能继续下滑，也能为后续问题二的珍稀物种响应提供上游输入。

因此，后文若出现底层资源暂时走低，也应结合时间滞后和局地压力一起解释，不能与修复结论直接对立。也只有把这一点写明，问题一的图和结论才不会只剩下单纯的数值对比，而是能形成“政策释放压力--食压重新分配--资源再平衡”这一完整链条。

\subsection{问题二分析}
问题二关心的是禁渔收益能否继续向珍稀物种传递，而不是单纯展示某个珍稀物种的数量是否增加。江豚和长江鲟都依赖食源恢复，但二者的约束条件并不相同：前者更接近对中层猎物供给的响应，后者则同时受到通道阻隔、栖息地破碎和人工放流的影响。若把它们都写成同一个增长率，模型就会把“自然恢复”和“补偿性增殖”混在一起，失去对机制差异的刻画。

因此，问题二更像是在搭建一个“珍稀物种响应模块”。题面附件中的事实需要逐一映射到参数：2022 年公报对应种群初值和监测基准，附件 1 中的航道阻隔、鱼道连通与栖息地破碎化对应 \(\delta_B\)，附件 2 中的放流、补充和食源恢复对应 \(R_S\)，而 \(e_D,m_D,H,u\) 则分别反映繁殖效率、自然死亡、捕捞/扰动和污染压力。只有把这些事实放进统一框架，才能说明禁渔收益是如何在食物网和生境连通性上被放大或削弱的，也才能解释为什么同处一个背景下，江豚和长江鲟会表现出不同的恢复曲线\cite{report2022}。

题面中的事实映射还意味着，问题二必须把“自然恢复”和“工程补偿”分开解释。对于江豚，更多是食源改善后的数量响应；对于长江鲟，除了食源，还要看洄游通道是否连续、放流是否提供额外补充、污染是否抵消恢复收益。因此，本问的重点不是比较哪个物种的绝对值更大，而是说明政策收益在不同约束条件下被保留了多少、损失了多少。把这层关系写清楚，才能让问题二从“珍稀物种是否增加”上升为“珍稀物种响应机制如何分化”。

表格的作用也在于让读者一眼看出每个参数到底来自哪条事实链，而不是把变量当成纯数学符号使用。这样写，问题二的证据链就会更完整：上游是禁渔和食源恢复，中间是通道、放流与污染修正，下游才是珍稀物种的数量变化。
